% TikZ 图：含电容器的闭合直流电路
% Day 20：闭合电路欧姆定律
\begin{tikzpicture}[scale=0.9,>=Latex]
  \tikzset{wire/.style={thick}, res/.style={thick}}
  \draw[wire] (0,0) -- (0,2.2) -- (0.9,2.2);
  \draw[res] (0.9,1.92) rectangle (1.8,2.48);
  \node[above] at (1.35,2.52) {$r$};
  \draw[wire] (1.8,2.2) -- (2.55,2.2);
  \draw[res] (2.55,1.92) rectangle (3.55,2.48);
  \node[above] at (3.05,2.52) {$R_1$};
  \draw[wire] (3.55,2.2) -- (4.4,2.2);
  \draw[res] (4.4,1.92) rectangle (5.45,2.48);
  \node[above] at (4.92,2.52) {$R_2$};
  \draw[wire] (5.45,2.2) -- (6.2,2.2) -- (6.2,0) -- (0,0);

  % capacitor in parallel with R2
  \draw[wire] (4.05,2.2) -- (4.05,1.08) -- (4.58,1.08);
  \draw[wire] (5.78,2.2) -- (5.78,1.08) -- (5.25,1.08);
  \draw[thick] (4.58,0.72) -- (4.58,1.44);
  \draw[thick] (5.25,0.72) -- (5.25,1.44);
  \node[below] at (4.92,0.68) {$C$};

  % source
  \draw[wire] (-0.35,0.62) -- (0.35,0.62);
  \draw[wire] (-0.50,1.25) -- (0.50,1.25);
  \node[left] at (-0.54,1.25) {$+$};
  \node[right] at (0.52,0.95) {$E$};
  \draw[->,red,thick] (2.02,2.52) -- (2.50,2.52) node[midway,above] {$I$};
  \node[font=\small,align=center] at (3.15,-0.45)
    {电容器与 $R_2$ 并联，稳定后 $U_C=U_{R_2}$};
\end{tikzpicture}
